\section{刑法论 - 刑法的原则}
\begin{enumerate}
  \item 罪刑法定。保障人权 > 保护法益。
        \begin{enumerate}
          \item 思想基础。
                \begin{enumerate}
                  \item 民主主义。
                  \item 自由主义。行动自由、\zd{国民的预测可能性}。刑不可知则威不可测。\al{流氓罪}。
                  \item 现存口袋罪。非法经营、寻衅滋事、聚众淫乱。
                \end{enumerate}
          \item 具体内容。
                \begin{enumerate}
                  \item 成文的罪刑法定。
                        \begin{enumerate}
                          \item 法律主义。只能全国人大定，国务院都没有。
                          \item 禁止习惯法。习惯不明确、不成文，容易扩张刑罚权。
                        \end{enumerate}
                  \item 事前的罪刑法定。涉及溯及力、时间效力问题。
                  \item 严格的罪刑法定。适用法条的时候，严格解释，不能类推解释。
                  \item 确定的罪刑法定。
                        \begin{enumerate}
                          \item 条文要规定得很明确。
                          \item 刑罚不确定 \yjt \zd{禁止绝对不定刑}；刑罚幅度不确定 \yjt \zd{禁止绝对不定期刑}。
                          \item \d{禁止处罚不当罚}。刑法谦抑性/补充性。刑法万能主义、重刑主义。\al{微博法官，死刑起步}。
                          \item 禁止不均衡残虐的刑罚。人道主义的要求，重刑思想。
                        \end{enumerate}
                \end{enumerate}
        \end{enumerate}
  \item 罪刑相适应原则。量刑的标准，刑罚的尺度。三项指标。
        \begin{enumerate}
          \item \zd{法益侵害性}/客观危害性。
          \item 可谴责性/\zd{主观恶性}。\sjt 已然发生的犯罪事实。报应刑/责任刑。
          \item 人身危害性/\zd{再犯可能性}。预防刑/目的刑。
        \end{enumerate}
\end{enumerate}
